\subsection{Henselian fields with an effective infinite residue field}
\label{C:infinite-residue-henselian-subsection}

\begin{theorem}[Proposed: exact arithmetic and finite quotient tables with separably closed residue]
\label{C:infinite-residue-henselian-arithmetic}
Let $r\ge1$, let $S_r$ be the explicit field of
Theorem~\ref{C:separably-closed-effective-arithmetic},
and let $k_0/S_r$ be a supplied finite field
extension.  Set
\[
 A=k_0[u]_{(u)},\qquad
 \mathcal O=A^h,\qquad H=\operatorname{Frac}(\mathcal O).
                                                               \tag{C.3.IR1}
\]
Then $H$ and all supplied finite towers $E/H$
have decidable exact finite names, equality,
arithmetic, tower validity and ordinary
univariate factorization, including inseparable
towers and polynomials.

Uniformly one computes their normalized
discrete valuations, uniformizers, finite
residue extensions of $k_0$ with exact
tables and residue maps, and an actual
$\mathcal O$-basis of the valuation ring
$\mathcal O_E$.
For every requested $N\ge1$ the quotient
$\mathcal O_E/\pi_E^N\mathcal O_E$ is
computed as a finite-dimensional
$k_0$-algebra, with its full table and
exact reduction map.
Every finite polynomial system over
that quotient has a uniform ordinary
solution decision.

All these fields are henselian and
defectless and their completion extensions
are regular.  No coefficient-field
section of $\mathcal O_E$ onto its
residue field is assumed or constructed.
\end{theorem}
\begin{proof}
\emph{The base field and its henselization.}
The arithmetic of $k_0$ is effective
by Theorem~\ref{C:separably-closed-effective-arithmetic}.
Elements of $k_0(u)$ are exact rational
functions, with $u$-adic valuation
computed by their first nonzero terms.
Thus $A$ has exact names, residue
field $k_0$, residue lifts by constants,
and Smith matrix-kernel algorithms.
The exact-access part of
Theorem~\ref{C:effective-coherent-henselization-WG}
uses no $C_1$ condition and supplies
exact pointed \'{e}tale names, equality,
residue maps and all kernels over
$\mathcal O=A^h$.
It is a discrete valuation ring,
with uniformizer $u$ and residue
field $k_0$, by Tag 07QL
\cite[Lemmas 15.46.1 and 15.46.11]
{C-StacksCoherentHenselizations}.

The fraction field $H$ consequently
has exact arithmetic and equality.
For a nonzero integral name, repeatedly
test its residue and divide by $u$
when the residue is zero.
The required quotient is found by
enumerating exact integral names,
and the process stops at its finite
valuation.  Fractions are treated
by subtracting valuations.
An integral fraction can be converted
to an integral name by these same
unit factorizations, or by exact
name enumeration once integrality
has been decided.
In particular all the reductions
\[
       \mathcal O/u^M\mathcal O
                  \simeq k_0[u]/(u^M)                \tag{C.3.IR2}
\]
are computed by repeated residue
extraction and division, using the
actual embedded constant field $k_0$.

The valued field $H$ is the
henselization of the one-variable
function field $k_0(u)$ at a
valuation trivial on $k_0$.
Jarden's Lemma~2.2(b) applies to
arbitrary constant fields, so it
proves defectlessness here, including
purely inseparable finite extensions
\cite[Lemmas 2.1--2.2]{C-JardenHenselian1994}.
It follows that for every finite
$E/H$ there is a unique extended
valuation, with
\[
 [E:H]=e(E/H)f(E/H).                                    \tag{C.3.IR3}
\]
Finite extensions remain henselian.
The degree, ramification-index and
residue-degree tower formulas show
that they also remain defectless.
Jarden's Lemma~2.1 therefore makes
$\widehat E/E$ regular for every
such $E$.  In particular its
separability hypothesis in positive
characteristic holds in Greenberg's
approximation theorem.

\emph{A complete residue-field certificate.}
Assume $E/H$ has been supplied as
a valid finite tower and put $n=[E:H]$.
All $H$-linear field operations and
the norm are computable in its
finite basis.
The unique extended valuation,
normalized to restrict to $v_H$,
is
\[
        w(x)=v_H(N_{E/H}(x))/n\quad(x\ne0),
                  \qquad w(0)=+\infty .               \tag{C.3.IR4}
\]
This follows by taking valuations
of all conjugates with their
inseparable multiplicities; it
is the same exact norm argument
as in
Theorem~\ref{C:henselian-global-field-arithmetic}.

Enumerate $e,f\ge1$ with $ef=n$,
finite field towers $T/k_0$ of
degree $f$, an explicit basis
$b_1=1,b_2,\ldots,b_f$ of $T$,
and elements $\pi,\beta_1,\ldots,\beta_f$
of $E$ with $\beta_1=1$.
The arithmetic theorem for $S_r$
decides validity of every field
tower $T/k_0$.
Write its computed multiplication
table as
\[
                 b_i b_j=\sum_\ell c_{ij\ell}b_\ell,
                       \qquad c_{ij\ell}\in k_0.
\]
Require
\[
 \begin{split}
 w(\pi)&=1/e,\qquad w(\beta_j)\ge0,\\
 w\left(\beta_i\beta_j-
                 \sum_\ell c_{ij\ell}\beta_\ell\right)&>0
                       \quad\text{for all }i,j.
 \end{split}                                           \tag{C.3.IR5}
\]
These are finite decidable tests.
The constants $k_0$ have valuation
zero away from zero and embed in
the actual residue field.
Thus (C.3.IR5) defines a unital
$k_0$-algebra homomorphism
$T\to\kappa_E$.
Since the source is a field,
it is injective.
If the actual ramification and
residue degrees are $e_0,f_0$,
the value of $\pi$ forces
$e\mid e_0$, and the residue
embedding forces $f\mid f_0$.
Equation (C.3.IR3) and $ef=n$
give $e=e_0$, $f=f_0$, so
the embedding is the full
residue field.

Such a certificate exists:
use an actual uniformizer,
the finite residue extension,
and arbitrary lifts of a
residue-field basis.
Every finite field extension
admits a finite tower
presentation, so it occurs
in the enumeration.
This remains valid when that
residue extension is purely
inseparable and nonmonogenic.
No residue primitive element
is required.
The search terminates and
computes $v_E=e w$, the
uniformizer $\pi_E=\pi$,
and the full field table $T$.

\emph{An actual integral basis.}
The successful certificate gives
the $n=ef$ elements
\[
         z_{ij}=\pi_E^i\beta_j,
             \qquad 0\le i<e,\quad1\le j\le f.         \tag{C.3.IR6}
\]
They form an $H$-basis of $E$
and an $\mathcal O$-basis of
$\mathcal O_E$.
Here is the precise valuation
argument, which also verifies
integrality without an assumed
coefficient section.
For coefficients $a_{ij}\in H$
and a fixed $i$, let
$m_i=\min_j v_H(a_{ij})$,
ignoring the all-zero block.
After dividing that block by
$u^{m_i}$, its coefficient
residues lie in $k_0$ and
are not all zero.
The residual $b_j$ are
$k_0$-linearly independent,
so no leading cancellation
occurs in $\sum_j a_{ij}\beta_j$.
Its value is exactly $m_i$.
Different $i$ have distinct
fractional parts $i/e$ modulo
$\mathbf Z$.  Consequently
\[
 w\left(\sum_{i,j}a_{ij}z_{ij}\right)
       =\min_{i,j}\bigl(v_H(a_{ij})+i/e\bigr)          \tag{C.3.IR7}
\]
for every nonzero combination.
This proves linear independence;
their number is $n$, so they
are a field basis.
Since $0\le i/e<1$ and the
coefficient valuations are
integers, (C.3.IR7) is nonnegative
exactly when every coefficient
is integral.  This proves
the asserted integral basis.

Finite $H$-linear algebra
computes the change of basis
from the supplied tower basis
to (C.3.IR6), and therefore
all multiplication constants.
The constants are integral
by (C.3.IR7) and are converted
to exact $\mathcal O$-names.
For an integral element in
this basis, its residue is
\[
           \sum_j\overline{a_{0j}}b_j\in T,            \tag{C.3.IR8}
\]
because the blocks with $i>0$
have positive value.
This is an exact residue
algorithm and provides lifts
by the $\beta_j$.

\emph{All finite quotient tables.}
Fix $N\ge1$ and choose $M$
with $eM\ge N$.
Then $u^M\mathcal O_E\subseteq
\pi_E^N\mathcal O_E$.
Reduce the integral multiplication
table in (C.3.IR6) modulo $u^M$
using (C.3.IR2).
This computes
$\mathcal O_E/u^M\mathcal O_E$
as an explicit finite-dimensional
$k_0$-algebra of dimension $nM$.
Compute the image of multiplication
by the known element $\pi_E^N$
in that algebra.  It is precisely
the ideal generated by $\pi_E^N$.
Quotienting this finite vector
space and its table gives
\[
              \mathcal O_E/\pi_E^N\mathcal O_E         \tag{C.3.IR9}
\]
with an exact $k_0$-basis,
multiplication table and
reduction map.
The actual constant field
$k_0\subset H\subset E$
is all that this construction
uses.  It does not lift the
larger residue field $T$
as a subfield of $E$.

Flatten (C.3.IR9) along the
known finite $S_r$-basis of
$k_0$ to get a finite-dimensional
$S_r$-algebra table.
Theorem~\ref{C:separably-closed-effective-arithmetic}
decides every finite polynomial
system and inequation over
that quotient by Haran's
algorithm on its coordinates.
The procedure is uniform in
$E,N$ and the input equations.

\emph{Factorization and the tower induction.}
Handle zero, nonzero constant,
and degree-one polynomials
directly and record the
nonzero leading scalar.
For a monic polynomial $f$
of degree $d\ge2$, a computed
substitution
\[
             F(Y)=\pi_E^{sd}f(Y/\pi_E^s)              \tag{C.3.IR10}
\]
with $s$ sufficiently large
makes all coefficients integral.
Every monic factor of $F$
over $E$ has integral
coefficients, because its
roots are integral and
$\mathcal O_E$ is integrally
closed.

For each possible proper
factor degree, coefficient
comparison in $F=GH$ is
a finite polynomial system
over $\mathcal O_E$ for
monic integral $G,H$.
One search enumerates exact
factor pairs over $E$.
The complementary search
tests these systems modulo
$\pi_E^N$, for increasing $N$,
using the quotient decisions
just constructed.
It reports irreducibility
when all proper degrees
have no solution at one
common finite level.
Greenberg's Theorem~1 and
Corollary~2 say that an
integral finite polynomial
system has a solution if
it has solutions at every
such level
\cite[pp. 59--62]{C-Greenberg1966}.
Its positive-characteristic
hypothesis holds because
$\widehat E/E$ is regular.
Thus if $F$ is irreducible,
each of the finitely many
factor systems fails at some
finite level, and a common
larger level makes the
negative search terminate.
If reducible, the positive
search terminates.
Recurse by degree and undo
(C.3.IR10) to obtain the
full factorization, including
inseparable factors.

Initially $E=H$ has the
explicit quotients (C.3.IR2),
so this factorization algorithm
already applies.
Before extending a valid
tower by a defining polynomial,
test its irreducibility over
the preceding field.
Once accepted, the new field
has a known finite $H$-basis,
so (C.3.IR4)--(C.3.IR9)
construct all its quotient
data directly over the fixed
base $H$ and $k_0$.
This inductive order proves
decidable tower validity and
all claimed algorithms without
circularity.
\end{proof}

\begin{theorem}[Proposed: complete ordinary split groups with infinite separably closed residue]
\label{C:infinite-residue-henselian-WG}
Let $H$ be as in (C.3.IR1), and let
$k/H$ be any supplied finite field tower.
Complete ordinary canonical classification,
actual isometry, equality, multiplication
and every specified exterior operation
for $WG_{\rm sp}(BO_{m,R})$, all $m\ge0$,
hold over:
\begin{enumerate}
\item every finite-dimensional commutative
$k$-algebra $R$ given by its multiplication
table;
\item every effective coherent local
henselization $R=B^h$ with the exact
name, arithmetic, residue-map and
finite matrix-kernel data of
Theorem~\ref{C:henselian-global-coefficient-WG},
whose residue field is a supplied
finite extension of $k$.
\end{enumerate}
The literal canonical description is
(C.3.HW3), with computed division
tables, adjoint involutions, and
decidable least-Gram-pair groups.
All finite coordinates have actual
positive and negative representatives.
Neither a division-index nor a
hermitian-isometry oracle is assumed.
\end{theorem}
\begin{proof}
We first verify the two arithmetic
inputs needed in the full proof
of Theorem~\ref{C:henselian-global-coefficient-WG}.
For every finite field extension
$F/H$, Theorem~\ref{C:infinite-residue-henselian-arithmetic}
gives exact arithmetic and factorization,
a henselian discrete valuation,
regular completion and all exact
finite quotient tables.
Haran's procedure decides every
polynomial system over those
quotients.  Therefore
Lemma~\ref{C:Greenberg-congruence-decisions}
decides rational points on every
explicit projective $F$-scheme
and constructs exact points
when they exist.
These are the proved Greenberg
and Haran procedures, not a
new hypothesis of existential
decidability over $F$ itself.

Equivalently,
Corollary~\ref{C:Greenberg-valued-algebra-decisions}
now applies over every finite $F/H$.
To make its role explicit, apply
this projective decision
to the explicit Grassmannians
of left ideals constructed in
Theorem~\ref{C:henselian-central-simple-algorithms}.
For a degree-$n$ central simple
algebra the least positive $j$
with a left ideal of dimension
$nj$ is its division index.
The finite list of projective
decisions computes that index;
the exact primitive-idempotent
and matrix-unit searches in
that proof then compute the
actual division algebra and
matrix decomposition.
This argument uses no finite
residue assumption beyond the
projective decision that has
just been independently supplied.
It handles every index in
the present characteristic.

The nilpotent-kernel certificate
of Theorem~\ref{C:henselian-global-finite-algebra-algorithms}
now computes the radical and
actual semisimple division
decomposition of any finite
algebra over $F$.
Every candidate center is a
finite tower over $F$, whose
validity and arithmetic are
decided by the preceding
arithmetic theorem; Artinian
structure supplies a true
certificate, including possible
inseparable centers.

For a division factor with
involution, compute its fixed
center by the linear kernel
of $\sigma-1$ on its center.
The adjoint fixes the coefficient
field, so that fixed center is
again finite over $H$.
Successive basis adjunctions
give an exact finite tower
presentation by linear algebra.
It therefore has the same
projective decision.
For two nonsingular hermitian
forms of equal rank, use the
explicit projective scheme of
half-dimensional invariant
isotropic subspaces of their
orthogonal difference, as in
Theorem~\ref{C:henselian-hermitian-isometry}.
Its points are exactly the
division-module Lagrangians.
Hyperbolicity and Witt cancellation
over the division algebra
equate their existence with
actual isometry.
This decides every residual
hermitian comparison, and
positive exact enumeration
constructs an isometry matrix.
The rank-zero case is immediate.
The same finite predecessor
tests give literal canonical
Gram classes and cancellative
Gram-pair group coordinates.

These arguments establish
exactly the finite-algebra
and hermitian inputs used in
the complete coefficient-ring
proof of
Theorem~\ref{C:henselian-global-coefficient-WG}.
All its remaining steps are
independent of the arithmetic
type of that field.
More explicitly, finite-algebra
projectivity is an affine
linear section test; actual
Hom modules and sign forms
are kernels; determinant
grids over the infinite field
decide invertibility; radical
idempotent lifts split actual
summands, with decreasing
dimension.
For $B^h$, finite free monic
neighborhoods and the supplied
matrix-kernel algorithm give
full actual kernels and
their actual residue images.
The endomorphism-image kernel
is radical, so its computed
finite residue quotient is
the actual semisimple quotient.
Lemma~\ref{C:nonnoetherian-henselian-lifting}
provides the exact idempotent,
form and isometry lifts.

Finite predecessor tests give
the literal indecomposable
indices and duality.
The sign-choice and reference
adjoint procedure gives the
computed division factors in
(C.3.HW3).
Residual coefficients lift
through actual generators,
are symmetrized by $1/2$,
and become actual nonsingular
forms because they are
invertible modulo the radical.
Thus the canonical Gram-pair
coordinates decide actual
isometry and group equality
and produce every finite
positive/negative representative.

Finally the finite tensor and
exterior constructions, determinant
pairings, and truncated inversion
of $\lambda_t$ in that proof
compute all the stated operations.
They involve finite projective
modules and invert no factorial.
This proves the complete ordinary
output for both coefficient cases.
\end{proof}

\begin{corollary}[Proposed: explicit non-$C_1$ fields with infinite residue]
\label{C:infinite-residue-new-families}
For every $r\ge1$ the field
\[
 H_r=\operatorname{Frac}
          \bigl(S_r[u]_{(u)}^h\bigr)                    \tag{C.3.IR11}
\]
has infinite separably closed residue
field $S_r$ and is not $C_1$.
It has transcendence degree $r+1$
over $\mathbf F_q$, so it is neither
a global function field nor a finite
extension of one of the henselized
global function fields with finite
residue treated earlier.
The preceding theorem therefore
gives a further whole family of
ordinary complete split-$WG$ answers,
including all its finite-dimensional
commutative coefficient algebras,
singular or nonreduced affine local
henselizations at supplied finite
residue extensions, and effective
coherent local henselizations with
those residue fields.
\end{corollary}
\begin{proof}
The $p$-basis of $S_r$ is
$t_1,\ldots,t_r$.
Adjoining the transcendental $u$
extends it to the $p$-basis
$t_1,\ldots,t_r,u$ of $S_r(u)$.
The henselization field is
separable algebraic over that
field, so the same basis
persists, by the finite-stage
degree argument in
Theorem~\ref{C:separably-closed-effective-arithmetic}.
Thus the homogeneous degree-$p$
form
\[
 \sum_{0\le a_1,\ldots,a_r,b<p}
            t_1^{a_1}\cdots t_r^{a_r}u^b
                      X_{a_1,\ldots,a_r,b}^p           \tag{C.3.IR12}
\]
in $p^{r+1}>p$ variables has
no nonzero $H_r$-zero.
This proves the non-$C_1$
assertion for every $r\ge1$.
The residue field is exactly
$S_r$ by henselization, and
the algebraic henselization
does not change transcendence
degree.

The finite coefficient algebras
are already covered by the
preceding theorem.
For finitely presented affine
local bases the usual exact
polynomial syzygies, localization
and pointed-neighborhood proof
of Lemma~\ref{C:C1-henselization-access}
give the required full kernel
algorithms over these effective
fields; no $C_1$ step is used
in that access argument.
This includes, for example,
$H_r[\epsilon]/(\epsilon^s)$
for $s\ge2$ and
$H_r[z]_{(z)}^h$.
All data and requested operations
are finite; no algorithm on
an infinitely supplied product
is asserted.
\end{proof}
